\section{Understand the code's complexity}
\label{section:preparation:labelling}

Software typically contains quite a lot of boilerplate code: 
predominantly administrative routines that enable calculations but don't deliver actual science flops. 
For several follow-on exercises, it will be very useful to know which code parts deliver science and which code parts do not directly add scientific value.
This is domain knowledge.
Using it contradicts, to some degree, our ambition to run a performance assessment in a black-box manner.


However, we do not ask an assessor to understand what a code does.
We require the assessor to know which code parts contribute towards the science case.
For this reason, it is absolutely sufficient to hand a first profile back to a domain expert and to ask them to label those routines that do the actual calculations. 

As an example: 
In a linear algebra code, the actual matrix-vector products, scalar products, \ldots~all deliver science.
The allocations for temporary block matrices, the exchange routines for rows of the matrices, any indexing over the sparsity pattern, and so forth do not contribute calculations.
They are absolutely essential, but they do not give us the MFlops/s.


\paragraph*{How it works}

\begin{itemize}[noitemsep]
  \item[$\boxplus$] Upon your very first profiling, ask a code expert which code parts (functions) actually perform calculations, i.e.~contribute towards the science.
\end{itemize}


\paragraph*{How it does not work}

\begin{itemize}[noitemsep]
  \item[$\boxminus$] Assume that all routines deliver science (even though you may assume that all are necessary).
\end{itemize}
